\documentclass[11pt,a4paper]{article}
\usepackage[T1]{fontenc}
\usepackage[utf8]{inputenc}
\usepackage[polish]{babel}
\usepackage[margin=2.3cm]{geometry}
\usepackage{amsmath,amssymb,booktabs,graphicx,array,longtable}
\usepackage[dvipsnames]{xcolor}
\usepackage[colorlinks=true,linkcolor=NavyBlue,urlcolor=NavyBlue]{hyperref}
\newcommand{\expv}{\langle O\rangle}
\setlength{\parskip}{3pt}

\title{QEM: wyniki publikowalne\\ oraz wyniki do uzupełnienia\\[4pt]
\large Zestawienie na podstawie danych z ODRY, Garneta i Emeralda}
\author{Karina Leśkiewicz}
\date{Październik 2026}

\begin{document}
\maketitle

\section*{Zakres i kryteria}

Zestawienie obejmuje wyłącznie wyniki, które (A) można opublikować w obecnej postaci, oraz
(B) wyniki, które staną się publikowalne po wskazanych uzupełnieniach. Wyniki, których nie
da się obronić bez zmiany tezy, pominięto (są omówione w recenzjach roboczych).

\textbf{Źródło danych:} repozytorium \url{https://github.com/karinaleskiewicz06/QEM},
gałąź \texttt{month2-topology-port}, commit \texttt{81e0f97}. Wszystkie liczby przeliczono
bezpośrednio z plików CSV wymienionych przy każdym wyniku.

\textbf{Kryterium A:} dane są w repozytorium, efekt jest istotny statystycznie albo
powtarza się w niezależnych kampaniach, a interpretacja nie wymaga założeń, których dane
nie testują. \textbf{Kryterium B:} kierunek efektu jest widoczny w danych, ale brakuje
kontroli, powtórzeń lub danych surowych; podano konkretną listę uzupełnień i kryterium
sukcesu.

\begin{table}[h]\centering\small
\caption{Przegląd wyników.}
\begin{tabular}{@{}lp{9.6cm}l@{}}
\toprule
& Wynik & Urządzenia \\
\midrule
\textbf{A1} & Twirling logiczny przy równej liczbie pomiarów obniża błąd CDR & ODRA, Garnet \\
\textbf{A2} & Błąd surowy rośnie z liczbą CZ; zysk ZNE jest największy przy najtańszym obwodzie & ODRA, Garnet \\
\textbf{A3} & Skończono-pomiarowa granica opłacalności ZNE ($\kappa$) dla trzech zbiorów węzłów & ODRA, Garnet \\
\midrule
B1 & Pełny (natywny) twirling, obejmujący CZ ze SWAP-ów & ODRA \\
B2 & Twirling zmniejsza wrażliwość CDR na dryf kalibracji & Garnet (ODRA: nie) \\
B3 & Kontrola z dopasowaną liczbą pomiarów jako wkład metodologiczny & symulacja \\
B4 & Degradację napędza liczba CZ i głębokość, nie sam routing & ODRA \\
B5 & Mapa koszt--dokładność: raw / ZNE / CDR / CDR+twirling & ODRA, Garnet \\
B6 & Górna granica okna ZNE (zanik sygnału przy dużej skali) & ODRA, Garnet, Emerald \\
B7 & Zgodność decyzji klasyfikatora z modelem idealnym po mitygacji & ODRA (Banknote), Garnet \\
B8 & Dynamiczne odsprzęganie jako wynik negatywny & ODRA \\
B9 & Nadwyżka spadku $r^2$ ponad szum niekoherentny & Garnet, ODRA \\
B10 & Mirror fidelity jako miara pomocnicza niezależna od zadania & ODRA, Garnet \\
B11 & Detekcja $[[4,2,2]]$ na obwodach Cliffordowskich & Garnet \\
B12 & Rozszerzenie wyników A1--A2 na Emerald & Emerald \\
\bottomrule
\end{tabular}
\end{table}

% =====================================================================
\section{Wyniki publikowalne}

\subsection*{A1. Twirling logiczny obniża błąd CDR na dwóch procesorach}

\textbf{Wynik.} Przy kontroli z dopasowaną liczbą pomiarów (ramię bez twirlingu uśrednia
$M$ identycznych kopii obwodu) twirling obniżył błąd CDR $\varepsilon$ w \textbf{16 z 17}
par sprzętowych (test znaków $p = 2{,}7\cdot10^{-4}$). W symulacji z szumem wyłącznie
niekoherentnym efektu nie ma, a w symulacji z wstrzykniętym obrotem koherentnym jest silny
(Tab.~\ref{tab:a1}, Rys.~\ref{fig:a1}).

\begin{table}[h]\centering\small
\caption{Pary ON/OFF dla CDR. Iloraz $=\varepsilon_{\text{ON}}/\varepsilon_{\text{OFF}}$.}
\label{tab:a1}
\begin{tabular}{@{}lcccccc@{}}
\toprule
Seria & okna & pary & $\varepsilon$ spadł & $p$ (znaki) & mediana ilorazu & średnie $\varepsilon$ OFF$\to$ON \\
\midrule
ODRA (sprzęt)        & 2 & 8  & 8/8   & 0{,}008 & 0{,}44 & 0{,}088 $\to$ 0{,}041 \\
Garnet (sprzęt)      & 3 & 9  & 8/9   & 0{,}039 & 0{,}66 & 0{,}055 $\to$ 0{,}033 \\
\midrule
sym.\ niekoherentna  & 5 & 15 & 10/15 & 0{,}30  & 0{,}92 & 0{,}051 $\to$ 0{,}045 \\
sym.\ koherentna $\theta=0{,}05$ & 5 & 15 & 15/15 & $6\cdot10^{-5}$ & 0{,}32 & 0{,}060 $\to$ 0{,}021 \\
\bottomrule
\end{tabular}
\end{table}

\textbf{Protokół.} ODRA: 8192 pomiary, $n_{\text{train}}=10$, $M=6$, 4 punkty ewaluacyjne,
drabina 12/30/39/57 CZ. Garnet: 1024 pomiary, $n_{\text{train}}=10$, $M=6$, 4 punkty,
drabina 12/57/96 CZ. Twirling na poziomie logicznym: bramki CZ pochodzące ze SWAP-ów nie są
owijane.

\textbf{Dane:} \texttt{odra\_twirl\_o1.csv}, \texttt{odra\_twirl\_o2.csv},
\texttt{garnet\_twirl\_w\{1,2,3\}.csv}, \texttt{night\_null\_w*.csv},
\texttt{night\_coh05\_w*.csv}.

\textbf{Co trzeba uczciwie napisać w pracy.} (i) Wariant natywny (pełne pokrycie CZ) nie
dał zysku i wymaga poprawki implementacji (B1); do tego czasu nie należy nazywać wariantu
logicznego ,,dolnym ograniczeniem'' pełnego randomized compiling. (ii) Zysk widać także przy
12~CZ, gdzie $r^2$ się nie zmienia, więc regułę decyzyjną trzeba opierać na $\varepsilon$,
a nie na $r^2$. (iii) Błąd liczony jest z 4 punktów ewaluacyjnych na punkt drabiny.

\begin{figure}[h]\centering
\includegraphics[width=\textwidth]{{{artifact:art_6a04e106-e401-4f02-97a1-46309fa21508}}}
\caption{(a) Błąd CDR z twirlingiem wobec błędu bez twirlingu, przy równej liczbie pomiarów;
każdy punkt to para (punkt drabiny $\times$ okno kalibracyjne); linia przerywana: $y=x$.
(b) Iloraz błędów dla każdej pary; kreska: mediana; nad seriami liczba par ze spadkiem błędu
i $p$ testu znaków.}
\label{fig:a1}
\end{figure}

\subsection*{A2. Błąd rośnie z liczbą CZ, a zysk ZNE jest największy przy najtańszym obwodzie}

\textbf{Wynik.} Na ODRZE błąd wartości oczekiwanej bez mitygacji rośnie monotonicznie z
liczbą CZ w dwóch niezależnych kampaniach, z różnymi modelami i różną liczbą punktów
(dokładne $p=0{,}083$ dla każdej, łącznie jednostronnie $p\approx0{,}002$). Na Garnecie
(8 punktów drabiny, rozmieszczenie zwarte i rozproszone) korelacja wynosi
$\rho_S=+0{,}93$ ($p=0{,}001$). Zysk ZNE (Richardson, skale $\{1,3,5\}$) jest największy przy
12~CZ (ODRA $4{,}6\times$ i $3{,}3\times$, Garnet $3{,}2\times$), a powyżej $\sim$30~CZ spada
do $1{,}0$--$1{,}7\times$ (Tab.~\ref{tab:a2}, Rys.~\ref{fig:a2}).

\begin{table}[h]\centering\small
\caption{Błąd surowy i zysk ZNE wzdłuż drabiny routingu.}
\label{tab:a2}
\begin{tabular}{@{}p{4.3cm}p{1.9cm}p{8.6cm}@{}}
\toprule
Kampania & CZ & błąd surowy / zysk ZNE \\
\midrule
ODRA 1 (model łatwy, 30 pkt, 4096, 3 powt.) & 12, 27, 36, 54 & 0{,}112 / 4{,}59; 0{,}281 / 2{,}44; 0{,}299 / 1{,}70; 0{,}357 / 1{,}49 \\
ODRA 2 (model trudny, 100 pkt, 8192) & 12, 27, 36, 54 & 0{,}075 / 3{,}30; 0{,}150 / 1{,}67; 0{,}186 / 1{,}32; 0{,}278 / 1{,}27 \\
Garnet (15 pkt, 4096, 1 okno) & 12--96 (8 pkt) & 0{,}061 / 3{,}17 przy 12 CZ; 0{,}075--0{,}205 / 1{,}03--1{,}34 \\
\bottomrule
\end{tabular}
\end{table}

\textbf{Dane:} \texttt{hardware\_results.csv} (ODRA 1, wiersze DD-off),
\texttt{odra\_full\_100pt.csv} (ODRA 2), \texttt{garnet\_full\_results.csv}.

\textbf{Co trzeba uczciwie napisać.} Na Garnecie spadek zysku powyżej 12~CZ nie jest
monotoniczny ($\rho_S(\text{CZ},\text{zysk})=-0{,}43$, $p=0{,}29$): wynik to ,,duży zysk
tylko przy najtańszym obwodzie'', a nie ,,monotoniczny spadek''. Przy 12~CZ błąd po ZNE leży
blisko podłogi szumu śrutowego, więc współczynnik $4{,}6\times$ jest ograniczony próbkowaniem.
Emerald nie pokazuje tego trendu (zysk 0{,}22 przy 21~CZ) i jest omówiony osobno (B12).

\begin{figure}[h]\centering
\includegraphics[width=\textwidth]{{{artifact:art_d240134b-5d51-4bcb-a956-86255c3064cc}}}
\caption{(a) Błąd $\expv$ bez mitygacji wobec liczby CZ po transpilacji. (b) Zysk ZNE
(błąd surowy / błąd po ZNE). Emerald pokazano dla kompletności; ma jedno okno i 4 punkty
ewaluacyjne.}
\label{fig:a2}
\end{figure}

\subsection*{A3. Granica opłacalności ZNE przy skończonym budżecie pomiarów}

\textbf{Wynik.} Dla ustalonego punktu drabiny różnica
$(\mathrm{MSE}_{\text{ZNE}}-\mathrm{MSE}_{\text{raw}})\cdot S$ jest liniowa w budżecie $S$,
a jej wyraz wolny wynosi $\kappa V$ ze stałą
$\kappa=|\Lambda|\sum_i c_i^2-1$, zależną tylko od węzłów ekstrapolacji. ZNE szkodzi poniżej
$S^{*}=\kappa V/B^2$, a pomaga powyżej. Zmierzone $\kappa$ zgadza się z przewidywaniem dla
trzech zbiorów węzłów na dwóch urządzeniach (Tab.~\ref{tab:a3}, Rys.~\ref{fig:a3}).
Niezależnie potwierdza się skalowanie $S^{*}\propto B^{-2}$: na Garnecie stosunek
$B^2$ dwóch punktów wynosi 3{,}07, a stosunek zmierzonych $S^{*}$ (847 wobec 291) 2{,}91.

\begin{table}[h]\centering\small
\caption{Stała $\kappa$: przewidziana i zmierzona (wyraz wolny regresji, 12 powtórzeń na
Garnecie, 5 na ODRZE).}
\label{tab:a3}
\begin{tabular}{@{}llcccc@{}}
\toprule
Węzły & Urządzenie, punkt & CZ & $\kappa$ przewidziane & $\kappa$ zmierzone & $S^{*}$ \\
\midrule
$\{1,3\}$ & Garnet, optimal:star & 12 & 4{,}00 & $3{,}54\pm0{,}56$ & 197 \\
$\{1,3\}$ & Garnet, adverse:star & 36 & 4{,}00 & $3{,}23\pm0{,}58$ & 76 \\
$\{1,3,5\}$ & Garnet, optimal:star & 12 & 14{,}66 & $14{,}42\pm0{,}95$ & 847 \\
$\{1,3,5\}$ & Garnet, adverse:star & 48 & 14{,}66 & $13{,}29\pm1{,}73$ & 291 \\
$\{1,3,5,7\}$ & ODRA, optimal:star & 12 & 44{,}56 & $48{,}1\pm6{,}0$ & 1160 \\
\bottomrule
\end{tabular}
\end{table}

\textbf{Dane:} \texttt{zne\_boundary\_k4\_results.csv}, \texttt{zne\_shot\_boundary\_results.csv},
\texttt{odra\_kappa\_k7.csv}.

\textbf{Jak to przedstawić.} Wyraz wolny wynika z propagacji wariancji dla niezależnych
pomiarów, a formalizm granicy pochodzi z pracy Scavino Alfaro (arXiv:2605.08251). Wkładem
jest więc \emph{eksperymentalna walidacja} tej granicy na sprzęcie i reguła praktyczna:
mniej węzłów obniża $S^{*}$ około czterokrotnie (z $\{1,3,5\}$ do $\{1,3\}$). Słowo
,,pre-rejestrowana'' należy usunąć albo podać commit z prognozą sprzed pomiaru. Punkt
ODRA all\_to\_all ($\kappa=41\pm49$) jest nieinformatywny i warto go tylko wymienić.

\begin{figure}[h]\centering
\includegraphics[width=\textwidth]{{{artifact:art_c8e610bd-ab2f-4eb2-9bcb-1bbeabf06dcd}}}
\caption{(a) Różnica MSE między ZNE a pomiarem bez mitygacji w funkcji budżetu pomiarów;
kółka: optimal:star, kwadraty: drugi punkt drabiny; linie: model $\kappa V/S + \text{const}$
z $\kappa$ przewidzianym z węzłów i nachyleniem z regresji. (b) $\kappa$ zmierzone wobec
przewidzianego; linia przerywana: $y=x$.}
\label{fig:a3}
\end{figure}

% =====================================================================
\section{Wyniki do uzupełnienia}

Każdy punkt ma ten sam układ: \emph{stan} (co pokazują obecne dane), \emph{brakuje},
\emph{do zrobienia}, \emph{kryterium} (co musi wyjść, żeby wynik był publikowalny).

\subsection*{B1. Pełny twirling (natywny)}
\textbf{Stan.} Owinięcie wszystkich CZ po transpilacji (\texttt{odra\_twirl\_native\_o1.csv})
obniżyło błąd tylko w 3 z 4 punktów (mediana ilorazu 0{,}93), a $r^2$ spadło średnio o 0{,}05.
\textbf{Brakuje.} Implementacja \texttt{twirl\_native} wstawia bramki Pauliego jako osobne
PRX bez scalania z sąsiednimi obrotami, przez co obwód ma 2{,}2--2{,}7$\times$ więcej bramek
jednokubitowych i około dwukrotnie większą głębokość niż obwód bez twirlingu.
\textbf{Do zrobienia.} Scalać bramki jednokubitowe po wstawieniu Pauli (np.\ przebieg
\texttt{Optimize1qGatesDecomposition} bez ponownego routingu) i sprawdzić, że liczba bramek
odpowiada obwodowi OFF; powtórzyć na ODRZE (2 okna) i Garnecie (3 okna) w tych samych
punktach co A1. \textbf{Kryterium.} Iloraz błędów nie gorszy niż w wariancie logicznym.

\subsection*{B2. Odporność CDR na dryf kalibracji}
\textbf{Stan.} Na Garnecie odchylenie błędu CDR między trzema oknami spada z
0{,}016 / 0{,}017 / 0{,}023 do 0{,}008 / 0{,}006 / 0{,}005 (12/57/96 CZ). Na ODRZE (2 okna)
efektu nie ma: przy 30 i 57~CZ rozrzut po twirlingu jest większy (Rys.~\ref{fig:b}a).
\textbf{Brakuje.} Liczby okien: dwa lub trzy okna nie pozwalają testować różnic wariancji.
\textbf{Do zrobienia.} Co najmniej 5 okien na urządzenie w tych samych punktach; test
Levene'a lub bootstrap ilorazu wariancji. \textbf{Kryterium.} Istotnie mniejsza wariancja
po twirlingu na obu urządzeniach.

\subsection*{B3. Kontrola z dopasowaną liczbą pomiarów}
\textbf{Stan.} Raport podaje, że bez dopasowania liczby pomiarów symulacja z szumem
depolaryzującym daje pozorny zysk $\Delta r^2\approx+0{,}10\ldots+0{,}15$ przy $M=10$, a po
dopasowaniu zysk znika. \textbf{Brakuje.} Pliku CSV z przebiegiem bez dopasowania w
repozytorium. \textbf{Do zrobienia.} Dołączyć dane symulacyjne i dodać jedno tanie ramię
sprzętowe bez dopasowania (np.\ Garnet, 2 punkty). \textbf{Kryterium.} Pozorny zysk widoczny
w symulacji i na sprzęcie, znikający po dopasowaniu.

\subsection*{B4. Liczba CZ i głębokość zamiast routingu}
\textbf{Stan.} Na ODRZE obwód bez routingu (gwiazda, $L=6/9/13$; \texttt{odra\_depth\_L*.csv})
traci jakość co najmniej tak szybko jak obwód z routingiem: $r^2$ 0{,}827 / 0{,}719 / 0{,}444
przy 24 / 36 / 52~CZ wobec 0{,}809 / 0{,}792 / 0{,}598 przy 27 / 36 / 54~CZ; błąd surowy przy
36~CZ 0{,}238 wobec 0{,}186 (Rys.~\ref{fig:b}b). \textbf{Brakuje.} Obwody bez routingu mają
więcej warstw, więc więcej bramek jednokubitowych i dłuższy czas trwania; jedna kampania,
jedno urządzenie. \textbf{Do zrobienia.} Powtórzyć na Garnecie; dla każdego obwodu zapisać
liczbę CZ, liczbę PRX i czas trwania; regresja błędu względem tych trzech zmiennych.
\textbf{Kryterium.} Rozdzielenie wkładu CZ, bramek 1q i czasu; teza ,,SWAP-y nie są
szczególnie szkodliwe przy tej samej liczbie CZ'' albo jej odrzucenie.

\subsection*{B5. Mapa koszt--dokładność metod}
\textbf{Stan.} W kampanii twirlingowej na ODRZE CDR z twirlingiem ma mniejszy błąd niż ZNE
w każdym punkcie (12 / 30 / 39 / 57~CZ: 0{,}022 / 0{,}041 / 0{,}061 / 0{,}040 wobec
0{,}026 / 0{,}067 / 0{,}165 / 0{,}123), ale zużywa około 22$\times$ więcej pomiarów.
Na Garnecie ZNE i CDR mierzono w osobnych kampaniach, których błędy surowe różnią się do
1{,}7$\times$. \textbf{Brakuje.} Porównania w jednej kampanii przy znormalizowanym koszcie.
\textbf{Do zrobienia.} Jedna kampania na urządzenie z czterema ramionami (raw, ZNE, CDR,
CDR+twirling), ten sam $n_{\text{train}}$, raportowanie błędu w funkcji efektywnej liczby
pomiarów. \textbf{Kryterium.} Front Pareto koszt--błąd z przedziałami niepewności z co
najmniej 2 okien.

\subsection*{B6. Górna granica okna ZNE}
\textbf{Stan.} Kryterium retencji $R=|\expv(5)|/|\expv(1)|$ zmierzone na ODRZE uznaje punkt
39~CZ za ,,zachowany'' ($R=0{,}81$ i $0{,}67$), choć $|\expv(1)|=0{,}04$, czyli sygnał jest
na podłodze już przy $\lambda=1$. \textbf{Brakuje.} Normalizacji do wartości idealnej i
wartości per skala na Garnecie i Emeraldzie. \textbf{Do zrobienia.} Retencja
$|\expv(1)|/|\expv_{\text{ideal}}|$ albo stosunek sygnału do błędu standardowego; test
przeskalowania Köstera--Mauerera
$\hat{\expv}_0/(\tfrac{15}{8}\expv(1))$ i kontrola ,,garbage folding''.
\textbf{Kryterium.} Próg zaniku zdefiniowany niezależnie od wyniku ZNE i potwierdzony na
2 urządzeniach.

\subsection*{B7. Zgodność decyzji klasyfikatora z modelem idealnym}
\textbf{Stan.} W danych z Miesiąca 1 (Banknote, ODRA) ZNE przybliża trafność do wartości
idealnej w 8 z 9 konfiguracji, a w 5 z 9 trafność surowa przekracza idealną. Na Garnecie
przy 96~CZ trafność surowa (0{,}78) przekracza bezszumową trafność testową modelu (0{,}767).
Trafność względem etykiet nie mierzy więc działania mitygacji. \textbf{Brakuje.} Danych per
punkt dla Banknote w repozytorium i trafności idealnej na tych samych punktach.
\textbf{Do zrobienia.} Metryka: odsetek zgodnych decyzji
$\mathrm{sign}\,\expv_{\text{mit}}=\mathrm{sign}\,\expv_{\text{ideal}}$, test McNemara na
punktach, co najmniej 100 punktów. \textbf{Kryterium.} Istotny wzrost zgodności po
mitygacji przy wysokim CZ.

\subsection*{B8. Dynamiczne odsprzęganie jako wynik negatywny}
\textbf{Stan.} Na ODRZE włączenie natywnego DD zwiększyło błąd we wszystkich 4 punktach
drabiny (\texttt{hardware\_results.csv}, np.\ 0{,}112$\to$0{,}159 przy 12~CZ) i obniżyło
mirror fidelity we wszystkich 13 konfiguracjach (Miesiąc 1). \textbf{Brakuje.} Kontroli
kolejności: DD-off zawsze wykonywano przed DD-on, a para dla wzorca circular została
rozdzielona na dwa dni (30.07 i 01.08, według znaczników czasu w ID jobów).
\textbf{Do zrobienia.} Przeplot ABAB w jednej sesji, obwód kontrolny z długą bezczynnością
(w którym DD powinno pomagać), co najmniej dwie sekwencje DD, dwa dni.
\textbf{Kryterium.} Kara za DD utrzymuje się przy przeplocie, a obwód kontrolny pokazuje,
w jakich warunkach DD pomaga.

\subsection*{B9. Nadwyżka spadku $r^2$ ponad szum niekoherentny}
\textbf{Stan.} $r^2$ spada z liczbą CZ na wszystkich urządzeniach, ale na Garnecie spadek
jest ilościowo taki sam jak w symulacji z szumem wyłącznie niekoherentnym (0{,}969$\to$0{,}827
wobec 0{,}985$\to$0{,}835--0{,}871; Rys.~\ref{fig:b}c). Sam spadek $r^2$ nie dowodzi więc
koherencji. \textbf{Brakuje.} Symulacji z kalibracją urządzenia (obecna kontrola używa
kalibracji zastępczej) i modeli łączących szum niekoherentny z koherentnym.
\textbf{Do zrobienia.} Null skalibrowany do Garneta i ODRY, modele z $\theta\in\{0{,}02;0{,}05\}$
na tym tle, ten sam $n_{\text{train}}$ i liczba pomiarów co na sprzęcie.
\textbf{Kryterium.} $r^2$ na sprzęcie istotnie poniżej nulla; dopiero wtedy można mówić o
$r^2$ jako wskaźniku koherencji.

\subsection*{B10. Mirror fidelity jako miara pomocnicza}
\textbf{Stan.} Na ODRZE mirror fidelity maleje monotonicznie z kosztem routingu
(0{,}549$\to$0{,}094, Miesiąc 1). Na Garnecie koreluje silniej z głębokością
($\rho_S=+0{,}79$, $p=0{,}021$, $n=8$) niż z liczbą CZ ($+0{,}48$).
\textbf{Brakuje.} Randomizacji obwodów lustrzanych (obwody $UU^\dagger$ bez randomizacji
częściowo znoszą błąd koherentny), członu SPAM, większej liczby wejść (teraz 3).
\textbf{Do zrobienia.} Randomized mirror circuits z pomiarem macierzy odczytu, $\geq10$ wejść.
\textbf{Kryterium.} Monotoniczna zależność od wybranej zmiennej kosztu z przedziałami
niepewności.

\subsection*{B11. Detekcja $[[4,2,2]]$ na obwodach Cliffordowskich}
\textbf{Stan.} Na Garnecie (\texttt{garnet\_qec422\_results.csv}) post-selekcja obniża błąd
we wszystkich 8 punktach, np.\ 0{,}418$\to$0{,}223 przy 40~CZ (akceptacja 0{,}75) i
0{,}863$\to$0{,}786 przy 131~CZ (akceptacja 0{,}56). \textbf{Brakuje.} Porównania przy
równej liczbie zaakceptowanych pomiarów; odchylenie między 3 ziarnami sięga 0{,}24.
\textbf{Do zrobienia.} $\geq10$ ziaren, błąd w funkcji efektywnego kosztu $S/p_{\text{akc}}$,
porównanie z ZNE na tym samym obwodzie. \textbf{Kryterium.} Przewaga detekcji przy równym
koszcie w określonym zakresie CZ. Wynik negatywny dla klasyfikatora (symulacja) opisać
oddzielnie.

\subsection*{B12. Emerald}
\textbf{Stan.} Jedno okno, 4 punkty ewaluacyjne, brak ramienia z twirlingiem i brak wartości
per skala ZNE. Błąd surowy rośnie z CZ (0{,}037$\to$0{,}101), ale przy 21~CZ ZNE pogarsza
wynik ponad czterokrotnie (0{,}168), czego model granicy z A3 nie tłumaczy co do wielkości.
\textbf{Do zrobienia.} Co najmniej 2 okna, $\geq8$ punktów, ramię CDR+twirling, zapis wartości
per skala; null można zasymulować na aktywnych kubitach obwodu (obwód logiczny ma 5 kubitów).
\textbf{Kryterium.} Wyniki A1 i A2 powtórzone na trzecim urządzeniu.

\begin{figure}[h]\centering
\includegraphics[width=\textwidth]{{{artifact:art_d52565f0-4897-4398-bd33-b9b33322b10e}}}
\caption{(a) Odchylenie standardowe błędu CDR między oknami kalibracyjnymi, bez i z
twirlingiem (B2). (b) ODRA: $r^2$ dopasowania CDR dla obwodu z routingiem ($L=3$) i bez
routingu (gwiazda, $L=6/9/13$) (B4). (c) Garnet: $r^2$ na sprzęcie (średnia $\pm$ SD z 3 okien)
i w symulacji z szumem wyłącznie niekoherentnym (5 okien) (B9).}
\label{fig:b}
\end{figure}

% =====================================================================
\section{Proponowany artykuł}

\textbf{Roboczy tytuł:} \emph{Shot-matched Pauli twirling improves Clifford data regression
under routing overhead on superconducting processors, and the finite-shot boundary of
zero-noise extrapolation.}

\textbf{Rdzeń (już teraz):} A1 jako wynik główny, A2 jako kontekst kosztu routingu, A3 jako
walidacja granicy ZNE. \textbf{Najbardziej wzmacniające uzupełnienia (w tej kolejności):}
B1 (pełny twirling), B5 (mapa koszt--dokładność), B2 (dryf), B4 (CZ vs routing).
Pozostałe punkty B mogą trafić do materiałów uzupełniających albo do osobnej, krótszej pracy
(B8 jako wynik negatywny dla użytkowników IQM).

\textbf{Wymagania wspólne przed wysłaniem.} Przypiąć wersję Qiskit i zapisywać transpilowane
obwody razem z ID jobów (liczby CZ dla tych samych obwodów różnią się między skryptami);
podawać liczbę okien i punktów przy każdym wyniku; cytować Köster i Mauerer
(arXiv:2607.09360) oraz Scavino Alfaro (arXiv:2605.08251).

\end{document}
